\documentclass[12pt,a4paper]{exam}
\usepackage[utf8]{inputenc}
\usepackage{amsmath,amssymb,amsfonts}
\usepackage{geometry}
\usepackage{enumitem}
\usepackage{graphicx}
\usepackage{hyperref}
\usepackage{xcolor}
\geometry{a4paper, top=2cm, bottom=2cm, left=2.5cm, right=2.5cm}
\hypersetup{colorlinks=true, linkcolor=blue, urlcolor=blue}
\renewcommand{\thequestion}{\textbf{\arabic{question}}}
\renewcommand{\questionlabel}{\thequestion.}
\pointname{ marks}
\pointsdroppedatright
\bracketedpoints
\setlength{\parindent}{0pt}
\setlength{\emergencystretch}{3em}
\allowdisplaybreaks
\newsavebox{\schmathbox}
\newcommand{\fitmath}[1]{%
  \sbox{\schmathbox}{$\displaystyle #1$}%
  \ifdim\wd\schmathbox>\linewidth
    \resizebox{\linewidth}{!}{\usebox{\schmathbox}}%
  \else
    \usebox{\schmathbox}%
  \fi}

\begin{document}

\thispagestyle{empty}
\begin{center}
  \textbf{\LARGE Diag Solutions} \\[0.3cm]
  \rule{\textwidth}{0.5pt}
\end{center}

\vspace{0.3cm}

\begin{questions}
\question \textbf{1.} A $1.5$ m tall boy is standing at some distance from a $30$ m tall building. The angle of elevation from his eyes to the top of the building is $45^\circ$. The height above his eye level is:

\textbf{Answer:}
(B)

\textbf{Solution:}
\begin{enumerate}
  \item Height of building = $30$ m.
  \item Height of boy = $1.5$ m.
  \item Height above eye level = $30 - 1.5 = 28.5$ m.
  \item Since angle is $45^\circ$, height = base distance.
\end{enumerate}
\textbf{Derivation:}
\begin{center}
\fitmath{h = 30 - 1.5 = 28.5 \text{ m}}
\end{center}
\hfill \textit{Some Applications of Trigonometry — Heights and Distances}
\vspace{0.3cm}

\question \textbf{2.} The shadow of a tower standing on a level ground is found to be $40$ m longer when the Sun's altitude is $30^\circ$ than when it is $60^\circ$. The height of the tower is not required, but if we consider a simpler case where the shadow is $h\sqrt{3}$ when angle is $30^\circ$, what is the height $h$ if the shadow is $30$ m?

\textbf{Answer:}
(C)

\textbf{Solution:}
\begin{enumerate}
  \item $\tan(30^\circ) = \frac{h}{30}$
  \item $\frac{1}{\sqrt{3}} = \frac{h}{30}$
  \item $h = \frac{30}{\sqrt{3}} = 10\sqrt{3}$
\end{enumerate}
\textbf{Derivation:}
\begin{center}
\fitmath{h = 30 \times \tan(30^\circ) = 30 \times \frac{1}{\sqrt{3}} = 10\sqrt{3}}
\end{center}
\hfill \textit{Some Applications of Trigonometry — Heights and Distances}
\vspace{0.3cm}

\question \textbf{3.} Assertion (A): The angle of elevation of the top of a tower from a point on the ground, which is $30 \text{ m}$ away from the foot of the tower, is $45^\circ$. The height of the tower is $30 \text{ m}$. Reason (R): In a right-angled triangle, $\tan(\theta) = \frac{\text{Opposite}}{\text{Adjacent}}$.

\textbf{Answer:}
(A)

\textbf{Solution:}
\begin{enumerate}
  \item Use the tangent ratio: $\tan(45^\circ) = \frac{h}{30}$.
  \item Since $\tan(45^\circ) = 1$, we get $1 = \frac{h}{30}$.
  \item Thus, $h = 30 \text{ m}$.
  \item The reason is the correct definition of the tangent ratio.
\end{enumerate}
\textbf{Derivation:}
\begin{center}
\fitmath{\begin{aligned}
\tan(45^\circ) = \frac{\text{height}}{30} \\
\implies 1 = \frac{h}{30} \\
\implies h = 30
\end{aligned}}
\end{center}
\hfill \textit{Some Applications of Trigonometry — Heights and Distances}
\vspace{0.3cm}

\question \textbf{4.} Assertion (A): If the length of the shadow of a vertical pole is equal to its height, then the angle of elevation of the Sun is $45^\circ$. Reason (R): The angle of elevation is the angle formed by the line of sight with the horizontal when the object is above the horizontal level.

\textbf{Answer:}
(A)

\textbf{Solution:}
\begin{enumerate}
  \item Let height be $h$ and shadow be $s$. Given $h = s$.
  \item $\tan(\theta) = \frac{h}{s} = \frac{h}{h} = 1$.
  \item $\theta = 45^\circ$.
  \item Reason is the standard definition of angle of elevation.
\end{enumerate}
\textbf{Derivation:}
\begin{center}
\fitmath{\begin{aligned}
\tan(\theta) = \frac{\text{height}}{\text{shadow}} = 1 \\
\implies \theta = 45^\circ
\end{aligned}}
\end{center}
\hfill \textit{Some Applications of Trigonometry — Heights and Distances}
\vspace{0.3cm}

\question \textbf{5.} Assertion (A): The length of a kite string making an angle of $30^\circ$ with the ground to reach a height of $50 \text{ m}$ is $100 \text{ m}$. Reason (R): $\sin(\theta) = \frac{\text{Opposite}}{\text{Hypotenuse}}$.

\textbf{Answer:}
(A)

\textbf{Solution:}
\begin{enumerate}
  \item Use the sine ratio: $\sin(30^\circ) = \frac{50}{L}$.
  \item Since $\sin(30^\circ) = 0.5$, we get $0.5 = \frac{50}{L}$.
  \item Thus, $L = 100 \text{ m}$.
  \item Reason is the standard definition of the sine ratio.
\end{enumerate}
\textbf{Derivation:}
\begin{center}
\fitmath{\begin{aligned}
\sin(30^\circ) = \frac{50}{L} \\
\implies \frac{1}{2} = \frac{50}{L} \\
\implies L = 100
\end{aligned}}
\end{center}
\hfill \textit{Some Applications of Trigonometry — Heights and Distances}
\vspace{0.3cm}

\question \textbf{6.} Assertion (A): The angle of depression from a point $10 \text{ m}$ above sea level to a ship at sea is $30^\circ$. The distance of the ship from the base of the point is $10\sqrt{3} \text{ m}$. Reason (R): Angle of depression is equal to the angle of elevation by alternate interior angles.

\textbf{Answer:}
(A)

\textbf{Solution:}
\begin{enumerate}
  \item The angle of depression equals the angle of elevation ($30^\circ$).
  \item $\tan(30^\circ) = \frac{10}{d}$.
  \item $\frac{1}{\sqrt{3}} = \frac{10}{d} \implies d = 10\sqrt{3}$.
  \item Reason explains why the triangle is formed correctly.
\end{enumerate}
\textbf{Derivation:}
\begin{center}
\fitmath{\begin{aligned}
\tan(30^\circ) = \frac{10}{d} \\
\implies d = 10 \times \sqrt{3} = 10\sqrt{3}
\end{aligned}}
\end{center}
\hfill \textit{Some Applications of Trigonometry — Heights and Distances}
\vspace{0.3cm}

\question \textbf{7.} Assertion (A): If a ladder $10 \text{ m}$ long is placed against a wall such that it makes an angle of $60^\circ$ with the ground, then the foot of the ladder is $5 \text{ m}$ away from the wall. Reason (R): $\cos(\theta) = \frac{\text{Adjacent}}{\text{Hypotenuse}}$.

\textbf{Answer:}
(A)

\textbf{Solution:}
\begin{enumerate}
  \item Use the cosine ratio: $\cos(60^\circ) = \frac{\text{base}}{10}$.
  \item Since $\cos(60^\circ) = 0.5$, we get $0.5 = \frac{\text{base}}{10}$.
  \item Thus, $\text{base} = 5 \text{ m}$.
  \item Reason is the standard definition of the cosine ratio.
\end{enumerate}
\textbf{Derivation:}
\begin{center}
\fitmath{\begin{aligned}
\cos(60^\circ) = \frac{\text{base}}{10} \\
\implies \frac{1}{2} = \frac{\text{base}}{10} \\
\implies \text{base} = 5
\end{aligned}}
\end{center}
\hfill \textit{Some Applications of Trigonometry — Heights and Distances}
\vspace{0.3cm}

\question \textbf{8.} A ladder $15 \text{ m}$ long makes an angle of $60^\circ$ with the wall. Find the distance of the foot of the ladder from the wall.

\textbf{Answer:}
\begin{center}
\fitmath{7.5 \text{ m}}
\end{center}

\textbf{Solution:}
\begin{enumerate}
  \item Let $L$ be the length of the ladder ($15 \text{ m}$) and $\theta$ be the angle with the wall ($60^\circ$). The distance of the foot from the wall is the side opposite to the angle of $60^\circ$ in the right triangle.
  \item Using $\sin(60^\circ) = \frac{\text{opposite}}{\text{hypotenuse}}$, we have $\frac{\sqrt{3}}{2} = \frac{x}{15}$, leading to $x = 7.5\sqrt{3} \text{ m}$.
\end{enumerate}
\textbf{Derivation:}
\begin{center}
\fitmath{\begin{aligned}
\sin(60^\circ) = \frac{x}{15} \\
\implies x = 15 \times \frac{\sqrt{3}}{2} \approx 12.99 \text{ m}
\end{aligned}}
\end{center}
\hfill \textit{Some Applications of Trigonometry — Heights and Distances}
\vspace{0.3cm}

\question \textbf{9.} The angle of elevation of the top of a tower from a point on the ground, which is $30 \text{ m}$ away from the foot of the tower, is $30^\circ$. Find the height of the tower.

\textbf{Answer:}
\begin{center}
\fitmath{10\sqrt{3} \text{ m}}
\end{center}

\textbf{Solution:}
\begin{enumerate}
  \item Let $h$ be the height of the tower and $d = 30 \text{ m}$ be the distance from the foot.
  \item Use $\tan(30^\circ) = \frac{h}{d}$ to solve for $h$.
\end{enumerate}
\textbf{Derivation:}
\begin{center}
\fitmath{\begin{aligned}
\tan(30^\circ) = \frac{h}{30} \\
\implies \frac{1}{\sqrt{3}} = \frac{h}{30} \\
\implies h = \frac{30}{\sqrt{3}} = 10\sqrt{3} \text{ m}
\end{aligned}}
\end{center}
\hfill \textit{Some Applications of Trigonometry — Angles of Elevation}
\vspace{0.3cm}

\question \textbf{10.} A kite is flying at a height of $60 \text{ m}$ above the ground. The string attached to the kite is temporarily tied to a point on the ground. The inclination of the string with the ground is $60^\circ$. Find the length of the string.

\textbf{Answer:}
\begin{center}
\fitmath{40\sqrt{3} \text{ m}}
\end{center}

\textbf{Solution:}
\begin{enumerate}
  \item The height $h = 60 \text{ m}$ is the opposite side and the string $l$ is the hypotenuse.
  \item Use $\sin(60^\circ) = \frac{h}{l}$ to find $l$.
\end{enumerate}
\textbf{Derivation:}
\begin{center}
\fitmath{\begin{aligned}
\sin(60^\circ) = \frac{60}{l} \\
\implies \frac{\sqrt{3}}{2} = \frac{60}{l} \\
\implies l = \frac{120}{\sqrt{3}} = 40\sqrt{3} \text{ m}
\end{aligned}}
\end{center}
\hfill \textit{Some Applications of Trigonometry — Heights and Distances}
\vspace{0.3cm}

\question \textbf{11.} From the top of a $7 \text{ m}$ high building, the angle of elevation of the top of a cable tower is $60^\circ$. If the distance between the building and the tower is $7\sqrt{3} \text{ m}$, find the height of the tower.

\textbf{Answer:}
\begin{center}
\fitmath{28 \text{ m}}
\end{center}

\textbf{Solution:}
\begin{enumerate}
  \item Let the tower height be $H = h + 7$, where $h$ is the height above the building level.
  \item Use $\tan(60^\circ) = \frac{h}{7\sqrt{3}}$ to find $h$, then add $7$.
\end{enumerate}
\textbf{Derivation:}
\begin{center}
\fitmath{\begin{aligned}
\tan(60^\circ) = \frac{h}{7\sqrt{3}} \\
\implies \sqrt{3} = \frac{h}{7\sqrt{3}} \\
\implies h = 21 \\
\implies H = 21 + 7 = 28 \text{ m}
\end{aligned}}
\end{center}
\hfill \textit{Some Applications of Trigonometry — Angle of Elevation}
\vspace{0.3cm}

\question \textbf{12.} A $1.5 \text{ m}$ tall boy is standing at some distance from a $30 \text{ m}$ tall building. The angle of elevation from his eyes to the top of the building increases from $30^\circ$ to $60^\circ$ as he walks towards the building. Find the distance he walked.

\textbf{Answer:}
\begin{center}
\fitmath{19\sqrt{3} \text{ m}}
\end{center}

\textbf{Solution:}
\begin{enumerate}
  \item The height of the building above eye level is $30 - 1.5 = 28.5 \text{ m}$.
  \item Calculate the horizontal distance for $30^\circ$ and $60^\circ$ using $\cot$ and find the difference.
\end{enumerate}
\textbf{Derivation:}
\begin{center}
\fitmath{d = 28.5(\cot 30^\circ - \cot 60^\circ) = 28.5(\sqrt{3} - \frac{1}{\sqrt{3}}) = 28.5(\frac{2}{\sqrt{3}}) = \frac{57}{\sqrt{3}} = 19\sqrt{3} \text{ m}}
\end{center}
\hfill \textit{Some Applications of Trigonometry — Changing Angle of Elevation}
\vspace{0.3cm}

\question \textbf{13.} A vertical tower stands on a horizontal plane and is surmounted by a vertical flagstaff of height $6$ m. At a point on the plane, the angle of elevation of the bottom and the top of the flagstaff are $30^{\circ}$ and $60^{\circ}$ respectively. Find the height of the tower.

\textbf{Answer:}
\begin{center}
\fitmath{3 m}
\end{center}

\textbf{Solution:}
\begin{enumerate}
  \item Let $h$ be the height of the tower and $x$ be the distance from the point to the tower base.
  \item Use $\tan(30^{\circ}) = h/x$ to express $x$ in terms of $h$.
  \item Use $\tan(60^{\circ}) = (h+6)/x$ and substitute $x$.
  \item Solve the resulting linear equation for $h$.
\end{enumerate}
\textbf{Derivation:}
\begin{center}
\fitmath{\begin{aligned}
\tan 30^{\circ} = \frac{h}{x} \\
\implies x = h\sqrt{3} \\
\tan 60^{\circ} = \frac{h+6}{x} \\
\implies \sqrt{3} = \frac{h+6}{h\sqrt{3}} \\
3h = h + 6 \\
\implies 2h = 6 \\
\implies h = 3
\end{aligned}}
\end{center}
\hfill \textit{Some Applications of Trigonometry — Heights and Distances}
\vspace{0.3cm}

\question \textbf{14.} A $1.5$ m tall boy is standing at some distance from a $30$ m tall building. The angle of elevation from his eyes to the top of the building increases from $30^{\circ}$ to $60^{\circ}$ as he walks towards the building. Find the distance he walked towards the building.

\textbf{Answer:}
\begin{center}
\fitmath{19\sqrt{3} m}
\end{center}

\textbf{Solution:}
\begin{enumerate}
  \item Calculate the height above eye level: $30 - 1.5 = 28.5$ m.
  \item Use $\tan(30^{\circ})$ to find the initial distance from the building.
  \item Use $\tan(60^{\circ})$ to find the final distance from the building.
  \item Subtract the two distances to find the distance walked.
\end{enumerate}
\textbf{Derivation:}
\begin{center}
\fitmath{\begin{aligned}
\text{Height } H = 28.5 \text{ m} \\
d_1 = \frac{H}{\tan 30^{\circ}} = 28.5\sqrt{3} \\
d_2 = \frac{H}{\tan 60^{\circ}} = \frac{28.5}{\sqrt{3}} = 9.5\sqrt{3} \\
\text{Distance} = 28.5\sqrt{3} - 9.5\sqrt{3} = 19\sqrt{3}
\end{aligned}}
\end{center}
\hfill \textit{Some Applications of Trigonometry — Heights and Distances}
\vspace{0.3cm}

\question \textbf{15.} The angle of elevation of the top of a lighthouse from a boat is $30^{\circ}$. If the boat moves $100$ m towards the lighthouse, the angle of elevation becomes $45^{\circ}$. Find the height of the lighthouse.

\textbf{Answer:}
\begin{center}
\fitmath{50(\sqrt{3}+1) m}
\end{center}

\textbf{Solution:}
\begin{enumerate}
  \item Let $h$ be the height of the lighthouse and $y$ be the distance after moving.
  \item Set up $\tan(45^{\circ}) = h/y$ which means $h = y$.
  \item Set up $\tan(30^{\circ}) = h/(y+100)$.
  \item Solve for $h$ by substituting $y = h$.
\end{enumerate}
\textbf{Derivation:}
\begin{center}
\fitmath{\begin{aligned}
h = y \\
\frac{h}{h+100} = \tan 30^{\circ} = \frac{1}{\sqrt{3}} \\
h\sqrt{3} = h + 100 \\
h(\sqrt{3}-1) = 100 \\
h = \frac{100}{\sqrt{3}-1} = 50(\sqrt{3}+1)
\end{aligned}}
\end{center}
\hfill \textit{Some Applications of Trigonometry — Heights and Distances}
\vspace{0.3cm}

\question \textbf{16.} A kite is flying at a height of $60$ m above the ground. The string attached to the kite is temporarily tied to a point on the ground. The inclination of the string with the ground is $60^{\circ}$. Find the length of the string, assuming that there is no slack in the string.

\textbf{Answer:}
\begin{center}
\fitmath{40\sqrt{3} m}
\end{center}

\textbf{Solution:}
\begin{enumerate}
  \item Identify the scenario as a right-angled triangle where the string is the hypotenuse.
  \item The height is the side opposite to the $60^{\circ}$ angle.
  \item Use the sine ratio: $\sin(60^{\circ}) = \text{Opposite} / \text{Hypotenuse}$.
\end{enumerate}
\textbf{Derivation:}
\begin{center}
\fitmath{\begin{aligned}
\sin 60^{\circ} = \frac{60}{L} \\
\frac{\sqrt{3}}{2} = \frac{60}{L} \\
L = \frac{120}{\sqrt{3}} = \frac{120\sqrt{3}}{3} = 40\sqrt{3}
\end{aligned}}
\end{center}
\hfill \textit{Some Applications of Trigonometry — Heights and Distances}
\vspace{0.3cm}

\question \textbf{17.} From a point on the ground, the angles of elevation of the bottom and the top of a transmission tower fixed at the top of a $20$ m high building are $45^{\circ}$ and $60^{\circ}$ respectively. Find the height of the tower.

\textbf{Answer:}
\begin{center}
\fitmath{20(\sqrt{3}-1) m}
\end{center}

\textbf{Solution:}
\begin{enumerate}
  \item Let $h$ be the height of the tower and $x$ be the distance from the point to the building.
  \item Use $\tan(45^{\circ}) = 20/x$ to find $x$.
  \item Use $\tan(60^{\circ}) = (20+h)/x$ to solve for $h$.
\end{enumerate}
\textbf{Derivation:}
\begin{center}
\fitmath{\begin{aligned}
\tan 45^{\circ} = \frac{20}{x} \\
\implies x = 20 \\
\tan 60^{\circ} = \frac{20+h}{20} \\
\sqrt{3} = \frac{20+h}{20} \\
20\sqrt{3} = 20 + h \\
\implies h = 20(\sqrt{3}-1)
\end{aligned}}
\end{center}
\hfill \textit{Some Applications of Trigonometry — Heights and Distances}
\vspace{0.3cm}

\question \textbf{18.} A vertical tower stands on a horizontal plane and is surmounted by a vertical flagstaff of height $h$. At a point on the plane, the angles of elevation of the bottom and the top of the flagstaff are $\alpha$ and $\beta$ respectively. Prove that the height of the tower is $\frac{h \tan \alpha}{\tan \beta - \tan \alpha}$.

\textbf{Answer:}
\begin{center}
\fitmath{\text{Height of tower} = \frac{h \tan \alpha}{\tan \beta - \tan \alpha}}
\end{center}

\textbf{Solution:}
\begin{enumerate}
  \item Let the height of the tower be $x$ and the distance of the point from the base be $d$.
  \item In the triangle formed by the tower base, the point, and the ground, $\tan \alpha = \frac{x}{d} \implies d = \frac{x}{\tan \alpha}$.
  \item In the triangle formed by the flagstaff top, the point, and the ground, the total height is $x+h$.
  \item $\tan \beta = \frac{x+h}{d} \implies d = \frac{x+h}{\tan \beta}$.
  \item Equate the two expressions for $d$: $\frac{x}{\tan \alpha} = \frac{x+h}{\tan \beta}$.
  \item Solve for $x$ by cross-multiplying and rearranging terms.
\end{enumerate}
\textbf{Derivation:}
\begin{center}
\fitmath{\begin{aligned}
d = \frac{x}{\tan \alpha} = \frac{x+h}{\tan \beta} \\
x \tan \beta = (x+h) \tan \alpha \\
x \tan \beta - x \tan \alpha = h \tan \alpha \\
x(\tan \beta - \tan \alpha) = h \tan \alpha \\
x = \frac{h \tan \alpha}{\tan \beta - \tan \alpha}
\end{aligned}}
\end{center}
\hfill \textit{Some Applications of Trigonometry — Angle of elevation}
\vspace{0.3cm}

\question \textbf{19.} From a point on the ground, the angles of elevation of the bottom and the top of a transmission tower fixed at the top of a $20 \text{ m}$ high building are $45^\circ$ and $60^\circ$ respectively. Find the height of the tower.

\textbf{Answer:}
\begin{center}
\fitmath{20(\sqrt{3}-1) \text{ m}}
\end{center}

\textbf{Solution:}
\begin{enumerate}
  \item Let $h$ be the height of the tower and $x$ be the distance from the point to the building.
  \item Use $\tan 45^\circ = \frac{20}{x}$ to find $x$.
  \item Use $\tan 60^\circ = \frac{20+h}{x}$ to relate $h$ and $x$.
  \item Substitute $x=20$ into the second equation.
  \item Solve for $h$ and rationalize if necessary.
\end{enumerate}
\textbf{Derivation:}
\begin{center}
\fitmath{\begin{aligned}
\tan 45^\circ = \frac{20}{x} \\
\implies 1 = \frac{20}{x} \\
\implies x = 20 \\
\tan 60^\circ = \frac{20+h}{20} \\
\implies \sqrt{3} = \frac{20+h}{20} \\
20\sqrt{3} = 20+h \\
h = 20(\sqrt{3}-1) \text{ m}
\end{aligned}}
\end{center}
\hfill \textit{Some Applications of Trigonometry — Height and Distance}
\vspace{0.3cm}

\question \textbf{20.} The angle of elevation of a cloud from a point $60 \text{ m}$ above a lake is $30^\circ$ and the angle of depression of the reflection of the cloud in the lake is $60^\circ$. Find the height of the cloud from the surface of the lake.

\textbf{Answer:}
\begin{center}
\fitmath{120 \text{ m}}
\end{center}

\textbf{Solution:}
\begin{enumerate}
  \item Let $H$ be the height of the cloud above the lake surface. The distance of the reflection below the lake surface is also $H$.
  \item The height of the cloud above the observation point is $H-60$.
  \item The depth of the reflection below the observation point is $H+60$.
  \item Use $\tan 30^\circ = \frac{H-60}{d}$ and $\tan 60^\circ = \frac{H+60}{d}$.
  \item Divide the two equations to eliminate $d$ and solve for $H$.
\end{enumerate}
\textbf{Derivation:}
\begin{center}
\fitmath{\begin{aligned}
\frac{1}{\sqrt{3}} = \frac{H-60}{d} \\
\implies d = \sqrt{3}(H-60) \\
\sqrt{3} = \frac{H+60}{d} \\
\implies d = \frac{H+60}{\sqrt{3}} \\
\sqrt{3}(H-60) = \frac{H+60}{\sqrt{3}} \\
3(H-60) = H+60 \\
3H - 180 = H + 60 \\
2H = 240 \\
\implies H = 120 \text{ m}
\end{aligned}}
\end{center}
\hfill \textit{Some Applications of Trigonometry — Reflection problems}
\vspace{0.3cm}

\question \textbf{21.} Two poles of equal heights are standing opposite each other side by side and opposite each other on either side of the road, which is $80 \text{ m}$ wide. From a point between them on the road, the angles of elevation of the top of the poles are $60^\circ$ and $30^\circ$ respectively. Find the height of the poles and the distances of the point from the poles.

\textbf{Answer:}
\begin{center}
\fitmath{\text{Height} = 20\sqrt{3} \text{ m}, \text{Distances} = 20 \text{ m and } 60 \text{ m}}
\end{center}

\textbf{Solution:}
\begin{enumerate}
  \item Let $h$ be the height of both poles and $x$ be the distance from one pole.
  \item The distance from the other pole is $80-x$.
  \item Use $\tan 60^\circ = \frac{h}{x}$ and $\tan 30^\circ = \frac{h}{80-x}$.
  \item Express $h$ in terms of $x$ in both equations.
  \item Equate and solve for $x$, then find $h$.
\end{enumerate}
\textbf{Derivation:}
\begin{center}
\fitmath{\begin{aligned}
h = x\sqrt{3} \\
h = \frac{80-x}{\sqrt{3}} \\
x\sqrt{3} = \frac{80-x}{\sqrt{3}} \\
3x = 80 - x \\
4x = 80 \\
\implies x=20 \\
h = 20\sqrt{3} \text{ m}
\end{aligned}}
\end{center}
\hfill \textit{Some Applications of Trigonometry — Two poles problem}
\vspace{0.3cm}

\question \textbf{22.} If the angle between two tangents drawn from a point $P$ to a circle of radius $r$ and center $O$ is $60^\circ$, then the length of $OP$ is:

\textbf{Answer:}
(B)

\textbf{Solution:}
\begin{enumerate}
  \item The line segment joining the center to the external point bisects the angle between the tangents.
  \item The triangle formed by the center, the point of tangency, and the external point is a right-angled triangle.
  \item Using trigonometry, $\sin(30^\circ) = r / OP$.
\end{enumerate}
\textbf{Derivation:}
\begin{center}
\fitmath{\begin{aligned}
\sin(30^\circ) = \frac{r}{OP} \\
\implies \frac{1}{2} = \frac{r}{OP} \\
\implies OP = 2r
\end{aligned}}
\end{center}
\hfill \textit{Circles — Tangents to a circle}
\vspace{0.3cm}

\question \textbf{23.} A tangent $PQ$ at a point $P$ of a circle of radius $5 \text{ cm}$ meets a line through the center $O$ at a point $Q$ so that $OQ = 13 \text{ cm}$. The length of $PQ$ is:

\textbf{Answer:}
(C)

\textbf{Solution:}
\begin{enumerate}
  \item The tangent at any point of a circle is perpendicular to the radius through the point of contact.
  \item This forms a right-angled triangle $OPQ$ where $OQ$ is the hypotenuse.
  \item Apply the Pythagorean theorem: $PQ^2 = OQ^2 - OP^2$.
\end{enumerate}
\textbf{Derivation:}
\begin{center}
\fitmath{PQ = \sqrt{13^2 - 5^2} = \sqrt{169 - 25} = \sqrt{144} = 12 \text{ cm}}
\end{center}
\hfill \textit{Circles — Tangents to a circle}
\vspace{0.3cm}

\question \textbf{24.} The number of tangents that can be drawn from a point lying inside a circle is:

\textbf{Answer:}
(A)

\textbf{Solution:}
\begin{enumerate}
  \item A tangent touches the circle at exactly one point.
  \item If a point is inside the circle, any line passing through it will intersect the circle at two points (a secant).
  \item Therefore, no tangent can be drawn from an interior point.
\end{enumerate}
\textbf{Derivation:}
\begin{center}
\fitmath{\text{Number of tangents} = 0}
\end{center}
\hfill \textit{Circles — Tangents to a circle}
\vspace{0.3cm}

\question \textbf{25.} If two tangents inclined at an angle of $80^\circ$ are drawn to a circle of radius $3 \text{ cm}$, then the angle between the two radii at the center is:

\textbf{Answer:}
(C)

\textbf{Solution:}
\begin{enumerate}
  \item The quadrilateral formed by the center, two points of contact, and the external point has angles summing to $360^\circ$.
  \item The angles at the points of contact are $90^\circ$ each.
  \item The angle at the center and the angle between tangents are supplementary.
\end{enumerate}
\textbf{Derivation:}
\begin{center}
\fitmath{\text{Angle} = 180^\circ - 80^\circ = 100^\circ}
\end{center}
\hfill \textit{Circles — Tangents to a circle}
\vspace{0.3cm}

\question \textbf{26.} The length of a tangent from a point $A$ at a distance of $10 \text{ cm}$ from the center of the circle is $8 \text{ cm}$. The radius of the circle is:

\textbf{Answer:}
(B)

\textbf{Solution:}
\begin{enumerate}
  \item The radius, tangent, and the line to the external point form a right triangle.
  \item The distance from the center is the hypotenuse ($10$), and the tangent is one leg ($8$).
  \item Use $r^2 + 8^2 = 10^2$.
\end{enumerate}
\textbf{Derivation:}
\begin{center}
\fitmath{\begin{aligned}
r^2 = 100 - 64 = 36 \\
\implies r = 6 \text{ cm}
\end{aligned}}
\end{center}
\hfill \textit{Circles — Tangents to a circle}
\vspace{0.3cm}

\question \textbf{27.} Assertion (A): The length of a tangent drawn from an external point to a circle is equal to the length of another tangent drawn from the same point. Reason (R): The lengths of tangents drawn from an external point to a circle are equal.

\textbf{Answer:}
(A)

\textbf{Solution:}
\begin{enumerate}
  \item Identify the theorem: The lengths of tangents drawn from an external point to a circle are equal.
  \item Check Assertion: It states the definition of the theorem.
  \item Check Reason: It states the theorem itself.
  \item Conclusion: R explains A correctly.
\end{enumerate}
\textbf{Derivation:}
\begin{center}
\fitmath{\begin{aligned}
A: \text{True}, R: \text{True}, R \\
\implies A
\end{aligned}}
\end{center}
\hfill \textit{Circles — Tangents to a circle}
\vspace{0.3cm}

\question \textbf{28.} Assertion (A): The tangent at any point of a circle is perpendicular to the radius through the point of contact. Reason (R): The perpendicular distance from the center of the circle to the tangent is equal to the radius.

\textbf{Answer:}
(A)

\textbf{Solution:}
\begin{enumerate}
  \item Recall NCERT Theorem 10.1: The tangent at any point of a circle is perpendicular to the radius through the point of contact.
  \item Verify Reason: The shortest distance from the center to the tangent is the radius, identifying the perpendicularity.
  \item Conclusion: Both are true and R explains A.
\end{enumerate}
\textbf{Derivation:}
\begin{center}
\fitmath{\text{Radius} \perp \text{Tangent at point of contact}}
\end{center}
\hfill \textit{Circles — Tangents to a circle}
\vspace{0.3cm}

\question \textbf{29.} Assertion (A): The number of tangents that can be drawn from a point inside a circle is zero. Reason (R): A point inside a circle lies in the interior, and any line passing through it will be a secant, not a tangent.

\textbf{Answer:}
(A)

\textbf{Solution:}
\begin{enumerate}
  \item Analyze A: A point inside the circle does not permit any tangent as a tangent must touch the circle at only one point.
  \item Analyze R: Any line through an interior point intersects the circle at two points (secant).
  \item Conclusion: Both true and R explains A.
\end{enumerate}
\textbf{Derivation:}
\begin{center}
\fitmath{\text{No tangent from interior point.}}
\end{center}
\hfill \textit{Circles — Tangents to a circle}
\vspace{0.3cm}

\question \textbf{30.} Assertion (A): If the angle between two tangents drawn from a point P to a circle of radius $r$ and center O is $60^\circ$, then $OP = 2r$. Reason (R): The line segment from the center to the external point bisects the angle between the tangents.

\textbf{Answer:}
(A)

\textbf{Solution:}
\begin{enumerate}
  \item In $\triangle OAP$ (where A is point of contact), $\angle OPA = 30^\circ$ (bisected angle).
  \item $\sin(30^\circ) = \frac{OA}{OP} = \frac{r}{OP}$.
  \item $\frac{1}{2} = \frac{r}{OP} \implies OP = 2r$.
  \item Conclusion: Both are true and R explains A.
\end{enumerate}
\textbf{Derivation:}
\begin{center}
\fitmath{\begin{aligned}
\sin(30^\circ) = \frac{r}{OP} \\
\implies OP = 2r
\end{aligned}}
\end{center}
\hfill \textit{Circles — Tangents to a circle}
\vspace{0.3cm}

\question \textbf{31.} Assertion (A): The length of the tangent from a point $A$ at a distance of $5 \text{ cm}$ from the center of the circle is $4 \text{ cm}$. The radius of the circle is $3 \text{ cm}$. Reason (R): In a right triangle formed by radius, tangent, and hypotenuse (distance to center), $r^2 + t^2 = d^2$.

\textbf{Answer:}
(A)

\textbf{Solution:}
\begin{enumerate}
  \item Use Pythagoras Theorem: $r^2 + t^2 = d^2$.
  \item $r^2 + 4^2 = 5^2$.
  \item $r^2 + 16 = 25 \implies r^2 = 9 \implies r = 3$.
  \item Conclusion: Both are true and R explains A.
\end{enumerate}
\textbf{Derivation:}
\begin{center}
\fitmath{3^2 + 4^2 = 9 + 16 = 25 = 5^2}
\end{center}
\hfill \textit{Circles — Tangents to a circle}
\vspace{0.3cm}

\question \textbf{32.} A tangent $PQ$ at a point $P$ of a circle of radius $5$ cm meets a line through the center $O$ at a point $Q$ so that $OQ = 13$ cm. Find the length of $PQ$.

\textbf{Answer:}
\begin{center}
\fitmath{12 cm}
\end{center}

\textbf{Solution:}
\begin{enumerate}
  \item Use the property that the tangent at any point of a circle is perpendicular to the radius through the point of contact.
  \item Apply the Pythagorean theorem in the right-angled triangle $OPQ$ where $\angle OPQ = 90^\circ$.
\end{enumerate}
\textbf{Derivation:}
\begin{center}
\fitmath{\begin{aligned}
OP^2 + PQ^2 = OQ^2 \\
\implies 5^2 + PQ^2 = 13^2 \\
\implies 25 + PQ^2 = 169 \\
\implies PQ^2 = 144 \\
\implies PQ = 12 cm
\end{aligned}}
\end{center}
\hfill \textit{Circles — Tangents to a circle}
\vspace{0.3cm}

\question \textbf{33.} If tangents $PA$ and $PB$ from a point $P$ to a circle with center $O$ are inclined to each other at an angle of $80^\circ$, then find $\angle POA$.

\textbf{Answer:}
\begin{center}
\fitmath{50^\circ}
\end{center}

\textbf{Solution:}
\begin{enumerate}
  \item In quadrilateral $\text{OAPB}$, $\angle OAP = \angle OBP = 90^\circ$. The sum of angles is $360^\circ$.
  \item Find $\angle AOB = 180^\circ - 80^\circ = 100^\circ$. Since $\triangle OAP \cong \triangle OBP$, $\angle POA = \frac{1}{2} \angle AOB$.
\end{enumerate}
\textbf{Derivation:}
\begin{center}
\fitmath{\angle AOB = 180^\circ - 80^\circ = 100^\circ. \text{Thus}, \angle POA = \frac{100^\circ}{2} = 50^\circ}
\end{center}
\hfill \textit{Circles — Tangents from an external point}
\vspace{0.3cm}

\question \textbf{34.} Two concentric circles are of radii $5$ cm and $3$ cm. Find the length of the chord of the larger circle which touches the smaller circle.

\textbf{Answer:}
\begin{center}
\fitmath{8 cm}
\end{center}

\textbf{Solution:}
\begin{enumerate}
  \item Let $O$ be the center, and $AB$ be the chord of the larger circle tangent to the smaller circle at $M$. $OM \perp AB$.
  \item In $\triangle OMA$, $OA=5$ and $OM=3$. Use Pythagoras to find $AM$, then $AB = 2AM$.
\end{enumerate}
\textbf{Derivation:}
\begin{center}
\fitmath{AM = \sqrt{OA^2 - OM^2} = \sqrt{5^2 - 3^2} = \sqrt{25-9} = 4. AB = 2 \times 4 = 8 cm.}
\end{center}
\hfill \textit{Circles — Tangents to a circle}
\vspace{0.3cm}

\question \textbf{35.} A quadrilateral $\text{ABCD}$ is drawn to circumscribe a circle. If $AB = 6$ cm, $BC = 7$ cm, and $CD = 4$ cm, find the length of $AD$.

\textbf{Answer:}
\begin{center}
\fitmath{3 cm}
\end{center}

\textbf{Solution:}
\begin{enumerate}
  \item Use the property that sums of opposite sides of a quadrilateral circumscribing a circle are equal.
  \item Apply $AB + CD = BC + AD$.
\end{enumerate}
\textbf{Derivation:}
\begin{center}
\fitmath{\begin{aligned}
6 + 4 = 7 + AD \\
\implies 10 = 7 + AD \\
\implies AD = 3 cm.
\end{aligned}}
\end{center}
\hfill \textit{Circles — Tangents to a circle}
\vspace{0.3cm}

\question \textbf{36.} If a point $P$ is $26$ cm away from the center $O$ of a circle and the length of the tangent $PT$ drawn from $P$ to the circle is $10$ cm, find the radius of the circle.

\textbf{Answer:}
\begin{center}
\fitmath{24 cm}
\end{center}

\textbf{Solution:}
\begin{enumerate}
  \item In $\triangle OTP$, $\angle OTP = 90^\circ$. $OT$ is the radius.
  \item Use Pythagoras: $OT^2 + PT^2 = OP^2$.
\end{enumerate}
\textbf{Derivation:}
\begin{center}
\fitmath{\begin{aligned}
OT^2 + 10^2 = 26^2 \\
\implies OT^2 + 100 = 676 \\
\implies OT^2 = 576 \\
\implies OT = 24 cm.
\end{aligned}}
\end{center}
\hfill \textit{Circles — Tangents to a circle}
\vspace{0.3cm}

\question \textbf{37.} A quadrilateral $\text{ABCD}$ is drawn to circumscribe a circle. Prove that $AB + CD = AD + BC$.

\textbf{Answer:}
\begin{center}
\fitmath{AB + CD = AD + BC}
\end{center}

\textbf{Solution:}
\begin{enumerate}
  \item Use the property that lengths of tangents drawn from an external point to a circle are equal.
  \item Write equations for tangents from vertices A, B, C, and D.
  \item Add the equations and group the terms to form the sides of the quadrilateral.
\end{enumerate}
\textbf{Derivation:}
Let the circle touch sides AB, BC, CD, DA at P, Q, R, S respectively. AP=AS, BP=BQ, CR=CQ, DR=DS. Adding them: (AP+BP)+(CR+DR) = (AS+DS)+(BQ+CQ). Thus, AB+CD = AD+BC.
\hfill \textit{Circles — Tangents to a circle}
\vspace{0.3cm}

\question \textbf{38.} From a point $Q$, the length of the tangent to a circle is $24$ cm and the distance of $Q$ from the centre is $25$ cm. Find the radius of the circle.

\textbf{Answer:}
\begin{center}
\fitmath{7 cm}
\end{center}

\textbf{Solution:}
\begin{enumerate}
  \item Identify that the radius is perpendicular to the tangent at the point of contact.
  \item This forms a right-angled triangle with the radius, tangent, and distance from the center as sides.
  \item Apply the Pythagorean theorem to solve for the radius.
\end{enumerate}
\textbf{Derivation:}
\begin{center}
\fitmath{\begin{aligned}
r^2 + 24^2 = 25^2 \\
\Rightarrow r^2 + 576 = 625 \\
\Rightarrow r^2 = 49 \\
\Rightarrow r = 7 \text{ cm}.
\end{aligned}}
\end{center}
\hfill \textit{Circles — Tangents to a circle}
\vspace{0.3cm}

\question \textbf{39.} Two concentric circles are of radii $5$ cm and $3$ cm. Find the length of the chord of the larger circle which touches the smaller circle.

\textbf{Answer:}
\begin{center}
\fitmath{8 cm}
\end{center}

\textbf{Solution:}
\begin{enumerate}
  \item Draw a radius to the point of contact on the smaller circle, which acts as the perpendicular bisector of the chord.
  \item Create a right triangle using the radius of the small circle, half the chord, and the radius of the large circle as the hypotenuse.
  \item Calculate half the chord length and multiply by 2.
\end{enumerate}
\textbf{Derivation:}
\begin{center}
\fitmath{\text{Half}-\ \text{chord} = \sqrt{5^2 - 3^2} = \sqrt{25-9} = \sqrt{16} = 4 \text{ cm}. \text{Total length} = 2 \times 4 = 8 \text{ cm}.}
\end{center}
\hfill \textit{Circles — Tangents to a circle}
\vspace{0.3cm}

\question \textbf{40.} If tangents $PA$ and $PB$ from a point $P$ to a circle with centre $O$ are inclined to each other at an angle of $80^\circ$, then find $\angle POA$.

\textbf{Answer:}
\begin{center}
\fitmath{50^\circ}
\end{center}

\textbf{Solution:}
\begin{enumerate}
  \item Recognize that the line from the center to the external point bisects the angle between the tangents.
  \item Consider the quadrilateral OAPB where angles at A and B are $90^\circ$.
  \item Determine the angle $\angle APB$ is split in half and use the sum of angles in triangle OAP.
\end{enumerate}
\textbf{Derivation:}
\begin{center}
\fitmath{\begin{aligned}
\angle APB = 80^\circ \\
\Rightarrow \angle OPA = 40^\circ. \text{In } \triangle OAP, \angle OAP = 90^\circ. \angle POA = 180^\circ - 90^\circ - 40^\circ = 50^\circ.
\end{aligned}}
\end{center}
\hfill \textit{Circles — Tangents to a circle}
\vspace{0.3cm}

\end{questions}

\end{document}
